\documentclass[12pt,a4paper]{article}

\usepackage[utf8]{inputenc}
\usepackage[T1]{fontenc}
\usepackage{lmodern}
\usepackage[english]{babel}
\usepackage{microtype}
\usepackage{amsmath}
\usepackage[margin=2.5cm]{geometry}
\usepackage{setspace}
\usepackage[hidelinks]{hyperref}
\usepackage{indentfirst}

\onehalfspacing
\setlength{\parindent}{1.25cm}
\setlength{\parskip}{0pt}
\setcounter{section}{2}

\begin{document}

\section{Results}
\label{sec:results}

\subsection{Sample and analytical coverage}
\label{sec-sample}

Of the fish entered into the study, 410 were retained for analysis: 136 in the sequential battery (BT1), 139 in the reversed behavioural context (BT2), and 135 in the isolated endurance context (BT3). Each treatment by exposure cell contained between 12 and 16 fish. Two endpoint-specific gaps remained after quality control: one fish in the sequential battery produced a no-choice first-choice record, and blood glucose was unavailable for 25 sequential-battery fish. The architecture of the three testing contexts is shown in Figure~1, and the composition of the treatment by exposure cells is given in Table~1.

\subsection{Primary measurement-context comparison}
\label{sec-context}

\subsubsection{Overall context pattern}
\label{sec-3-2-1}

Endpoints were compared between the sequential battery and the corresponding comparison context within matched treatment by exposure cells. Thirteen behavioural and performance endpoints were eligible, each contributing nine treatment by exposure cells, giving 117 comparisons in total. A Benjamini--Hochberg-adjusted context difference was detected in 19 of them, while 98 of 117 (83.8\%) showed no detected adjusted difference.

The detections were not evenly distributed. Among the light-dark and novel tank endpoints, compared against the reversed behavioural context, 8 of 81 comparisons showed a detected difference and 73 did not; among the endurance endpoints, compared against the isolated endurance context, 11 of 36 showed a detected difference and 25 did not. Where a difference was detected it was generally large, with a median absolute Hedges' $g$ of 1.16 across the 19 comparisons. Equivalence testing was not conducted because endpoint-specific equivalence margins were not prespecified.

Throughout this section, effect sizes are oriented as the comparison context minus the sequential battery, so a positive value indicates higher measurements in the comparison context and a negative value indicates lower measurements there. Figure~2 displays representative endpoints only; the complete set of estimates, bootstrap confidence intervals, and adjusted p-values for all 117 comparisons is given in Table~3.

\subsubsection{Light-dark and novel-tank contexts}
\label{sec-3-2-2}

Within the light-dark assay, 2 of 45 comparisons showed a detected adjusted difference, and both occurred in control fish at the 24 h exposure. Time in the bright zone was higher in the reversed behavioural context than in the sequential battery ($g$ = 1.10, 95\% CI [0.51, 1.88], $p_{\mathrm{BH}}$ = 0.041), with medians of 203.50 s in the reversed context against 97.00 s in the sequential battery (Figure~2A). Mean time per crossing followed the same direction ($g$ = 0.64, 95\% CI [0.25, 1.05], $p_{\mathrm{BH}}$ = 0.050; medians 10.36 s and 4.84 s respectively). No adjusted difference was detected in any first-choice, latency, or zone-transition comparison, or in the remaining time-in-bright-zone and mean-time-per-crossing cells (Table~3).

Within the novel tank assay, 6 of 36 comparisons showed a detected adjusted difference. Total distance was higher in the reversed behavioural context in control fish at 1 h ($g$ = 1.14, 95\% CI [0.44, 2.12], $p_{\mathrm{BH}}$ = 0.045; medians 6544.25 cm and 4258.30 cm respectively; Figure~2B). Mean velocity was higher in the reversed context in two cells: control fish at 1 h ($g$ = 1.10, 95\% CI [0.41, 2.20], $p_{\mathrm{BH}}$ = 0.041) and 0.5\% fish at 96 h ($g$ = 1.15, 95\% CI [0.56, 1.89], $p_{\mathrm{BH}}$ = 0.030). Total movement time ran the other way in two cells, being lower in the reversed context. In 1\% fish at 1 h the estimate was clear ($g$ = $-$0.75, 95\% CI [$-$1.89, $-$0.10], $p_{\mathrm{BH}}$ = 0.025; medians 295.50 s and 438.50 s). In control fish at 24 h the two summaries did not agree: the adjusted rank test detected a difference ($p_{\mathrm{BH}}$ = 0.028), while the bootstrap interval for Hedges' $g$ included zero ($g$ = $-$0.58, 95\% CI [$-$1.64, 0.11]). Both are reported as obtained.

The largest novel tank difference was in upper-stratum preference in 0.5\% fish at 24 h, which was higher in the reversed context ($g$ = 1.94, 95\% CI [0.93, 4.37], $p_{\mathrm{BH}}$ = 0.006; medians 91.77\% and 44.79\%). Total distance, mean velocity, and total movement time are related locomotor measurements. The remaining novel tank comparisons showed no detected adjusted difference (Table~3).

\subsubsection{Swimming-resistance context}
\label{sec-3-2-3}

The endurance endpoints, compared between the sequential battery and the isolated endurance context, accounted for 11 of the 19 detected differences and were distributed across all four endpoints.

The swimming-resistance index differed in four cells, and in every case the values in the isolated context were lower than those in the sequential battery: 0.5\% at 1 h ($g$ = $-$1.38, 95\% CI [$-$2.55, $-$0.66], $p_{\mathrm{BH}}$ = 0.018; medians 17.85 and 35.38), 1\% at 1 h ($g$ = $-$1.53, 95\% CI [$-$2.96, $-$0.84], $p_{\mathrm{BH}}$ = 0.023; medians 5.15 and 32.87), 0.5\% at 24 h ($g$ = $-$1.68, 95\% CI [$-$4.09, $-$0.84], $p_{\mathrm{BH}}$ = 0.017; medians 16.65 and 42.30), and 0.5\% at 96 h ($g$ = $-$1.16, 95\% CI [$-$2.32, $-$0.47], $p_{\mathrm{BH}}$ = 0.041; medians 8.17 and 32.20). Figure~2C shows the swimming-resistance index across contexts.

The final flow step followed the same direction in four cells, again lower in the isolated context: 0.5\% at 1 h ($g$ = $-$1.30, 95\% CI [$-$2.42, $-$0.61], $p_{\mathrm{BH}}$ = 0.018), 1\% at 1 h ($g$ = $-$1.36, 95\% CI [$-$2.63, $-$0.65], $p_{\mathrm{BH}}$ = 0.023), 0.5\% at 24 h ($g$ = $-$1.23, 95\% CI [$-$2.97, $-$0.47], $p_{\mathrm{BH}}$ = 0.017), and 0.5\% at 96 h ($g$ = $-$0.97, 95\% CI [$-$2.01, $-$0.27], $p_{\mathrm{BH}}$ = 0.045). Time at the final flow step differed in a single cell, 0.5\% at 24 h, and was also lower in the isolated context ($g$ = $-$1.34, 95\% CI [$-$2.59, $-$0.65], $p_{\mathrm{BH}}$ = 0.025; medians 16.50 s and 42.00 s).

The number of attempts ran in the opposite direction, being higher in the isolated context in two cells: 1\% at 1 h ($g$ = 1.37, 95\% CI [0.56, 2.66], $p_{\mathrm{BH}}$ = 0.017; medians 3.00 and 1.00) and 0.5\% at 96 h ($g$ = 0.89, 95\% CI [0.21, 1.78], $p_{\mathrm{BH}}$ = 0.026; medians 2.00 and 1.00). The complete endurance results, including the 25 comparisons in which no adjusted difference was detected, are given in Table~3.

\subsection{Exploratory relationships among endpoints}
\label{sec-explore}

\subsubsection{Correlation structure}
\label{sec:res-corr}

Rank correlations among the 15 registered endpoints in the sequential battery were computed on pairwise-complete observations, with pairwise sample sizes ranging from 110 to 136 (Figure~3). In the heatmap, a diagonal slash marks pairs whose correlation did not reach the unadjusted display criterion of p < 0.05.

The strongest associations lay within assay domains, and several of them reflect how the measurements are constructed. In the endurance assay the swimming-resistance index and the final flow step were almost completely collinear ($\rho$ = 0.98). Within the novel tank, total distance and mean velocity were strongly associated ($\rho$ = 0.87), whereas mean velocity and total movement time were essentially unrelated ($\rho$ = 0.05). Within the light-dark assay, time in the bright zone and mean time per crossing were closely associated ($\rho$ = 0.77), and first choice was associated with both ($\rho$ = 0.44 and 0.45 respectively).

Associations across domains were weaker but not absent. Light-dark zone transitions were associated with total distance ($\rho$ = 0.30) and total movement time ($\rho$ = 0.38), and upper-stratum preference was associated with the other novel tank measurements ($\rho$ = 0.44 with total distance, 0.37 with total movement time, and 0.34 with mean velocity). The number of attempts was inversely related to both endurance measurements, the swimming-resistance index ($\rho$ = $-$0.37) and the final flow step ($\rho$ = $-$0.36). Two relationships were weak: condition factor with the swimming-resistance index ($\rho$ = 0.14), and time in the bright zone with upper-stratum preference ($\rho$ = 0.09).

Endpoints that are operationally or mathematically related, such as the swimming-resistance index and the final flow step, or total distance and mean velocity, carry overlapping rather than independent information.

\subsubsection{Principal component analysis}
\label{sec-3-3-2}

The multivariate structure of the sequential battery was summarised with a single principal component analysis of 14 registered endpoints in complete cases ($n = 135$ of 136), with all 14 loading vectors displayed (Figure~4). Blood glucose was excluded from this analysis because requiring it as a complete case would have reduced the available sample from 135 to 110 fish; it is therefore retained in the 15-endpoint correlation analysis and in its endpoint-level model, but not here.

The first component accounted for 23.3\% of the variance and the second for 17.2\%, together 40.6\%. The third accounted for 10.3\% and appears only within the complete variance spectrum (Figure~S1). The first component was oriented most strongly towards total distance, mean velocity, total movement time, the swimming-resistance index, and the final flow step. Along the second component, bright-zone occupancy and the novel tank locomotor measurements pointed in the opposite direction to the swimming-resistance index and the final flow step. Component signs are arbitrary, and no separation between treatment or exposure groups is inferred from this display.

\subsection{Secondary ethanol responsiveness}
\label{sec-ethanol}

\subsubsection{Overall endpoint-level pattern}
\label{sec-3-4-1}

Ethanol effects were evaluated across 14 modelled endpoints within the sequential battery. Ten of these showed at least one Benjamini--Hochberg-adjusted treatment, exposure, interaction, or coefficient-level term under their governing model. Four showed no such term: latency under its hurdle model, the number of zone transitions, the number of attempts, and time at the final flow step. Condition factor was recorded as a descriptive control measure and was not among the 14 modelled endpoints.

For endpoints analysed by aligned-rank-transform ANOVA, inference is rank-based, and compact letter displays localise Benjamini--Hochberg-adjusted contrasts on that rank scale. Raw medians are therefore reported separately, in their original units.

\subsubsection{Light-dark test}
\label{sec-3-4-2}

Time in the bright zone showed treatment ($p_{\mathrm{BH}}$ = 1.23 $\times$ $10^{-5}$) and exposure ($p_{\mathrm{BH}}$ = 1.02 $\times$ $10^{-4}$) effects, with no detected interaction ($p_{\mathrm{BH}}$ = 0.090). At every exposure duration the 1\% group was distinguished from the control group, with the 0.5\% group intermediate and separated from neither. Raw medians rose with both concentration and exposure duration, reaching 439.00 s in the 1\% group at 96 h against 117.50 s in the corresponding control (Figure~5A; Tables~4 and 5).

Mean time per crossing showed treatment ($p_{\mathrm{BH}}$ = 7.51 $\times$ $10^{-5}$), exposure ($p_{\mathrm{BH}}$ = 2.25 $\times$ $10^{-5}$), and interaction ($p_{\mathrm{BH}}$ = 5.32 $\times$ $10^{-5}$) terms, and its localisation varied with exposure duration. At 1 h the interaction-aligned contrasts distinguished all three groups from one another. At 24 h the 1\% group differed from both the control and 0.5\% groups. At 96 h the 0.5\% group differed from both the control and 1\% groups, which were not distinguished from each other. These letters describe contrasts on the aligned-rank scale and are reported alongside, not in place of, the raw medians (Figure~5B; Tables~4 and 5).

First choice, analysed by Firth penalised logistic regression, showed no detected main effect, but four coefficient-level treatment by exposure terms were detected: 1\% at 24 h ($p_{\mathrm{BH}}$ = 0.002), 0.5\% at 96 h ($p_{\mathrm{BH}}$ = 0.003), 1\% at 96 h ($p_{\mathrm{BH}}$ = 0.017), and 0.5\% at 24 h ($p_{\mathrm{BH}}$ = 0.029) (Figure~S2A; Table~4).

Latency was not detected in either component of its governing two-part hurdle model, with all terms exceeding p > 0.23 in the binary component and p > 0.31 in the conditional component (Figure~S2B; Table~4). The number of zone transitions was not detected under a negative-binomial model (Figure~S3).

\subsubsection{Novel tank test}
\label{sec-3-4-3}

Total distance showed a treatment by exposure interaction ($p_{\mathrm{BH}}$ = 0.00117) with no detected main effects. The 1\% group was distinguished from the control group at 1 h and again at 96 h, but not at 24 h, and the raw medians ran in opposite directions at those two durations: 6668.30 cm against 4258.30 cm at 1 h, and 3911.70 cm against 5390.47 cm at 96 h (Figure~6A; Tables~4 and 5).

Mean velocity showed the same interaction structure ($p_{\mathrm{BH}}$ = 0.00644). At 1 h the 1\% group was distinguished from the control group with the 0.5\% group intermediate; no groups were distinguished at 24 h; and at 96 h the 1\% group was again distinguished from the control (Figure~S4). Total movement time also showed an interaction ($p_{\mathrm{BH}}$ = 0.029), localised entirely to the 24 h exposure, where the 1\% group was distinguished from both the control and 0.5\% groups and no groups were distinguished at 1 h or 96 h (Figure~S5).

Upper-stratum preference showed a treatment term ($p_{\mathrm{BH}}$ = 0.032), with exposure and the interaction not detected (both $p_{\mathrm{BH}}$ = 0.200). No pairwise treatment contrast reached significance after adjustment, so no localised group difference is reported (Figure~6B; Table~4).

\subsubsection{Swimming-resistance assay}
\label{sec-3-4-4}

The swimming-resistance index showed treatment ($p_{\mathrm{BH}}$ = 1.23 $\times$ $10^{-5}$) and exposure ($p_{\mathrm{BH}}$ = 0.029) effects, while the interaction did not reach the adjusted threshold ($p_{\mathrm{BH}}$ = 0.054). At each exposure duration the 1\% group was distinguished from both the control and 0.5\% groups. Raw medians in the 1\% group were 7.40 at 24 h and 9.67 at 96 h, against control medians of 44.25 and 36.09 respectively (Figure~7; Tables~4 and 5).

The final flow step, analysed by proportional-odds regression, showed a treatment term (likelihood-ratio $\chi^{2}(2)$ = 23.71, $p_{\mathrm{BH}}$ = 2.13 $\times$ $10^{-5}$) and a treatment by exposure term ($\chi^{2}(4)$ = 12.10, $p_{\mathrm{BH}}$ = 0.025); exposure was not detected ($\chi^{2}(2)$ = 4.95, $p_{\mathrm{BH}}$ = 0.084). This endpoint is reported in Table~4 and is not plotted separately. Because the final flow step and the swimming-resistance index are almost completely collinear (Section~\ref{sec:res-corr}), they are reported as related rather than independent evidence.

The number of attempts was not detected under a Poisson model (Figure~S6), and time at the final flow step was not detected under aligned-rank-transform ANOVA (Table~4).

\subsubsection{Blood glucose and condition factor}
\label{sec-3-4-5}

Blood glucose showed an exposure effect ($p_{\mathrm{BH}}$ = 0.022) and a treatment by exposure interaction ($p_{\mathrm{BH}}$ = 0.0019), with no detected treatment main effect ($p_{\mathrm{BH}}$ = 0.669). Localisation varied with exposure duration. At 1 h no pairwise difference was detected. At 24 h the 0.5\% and 1\% groups were distinguished from each other, with the control group intermediate and separated from neither. At 96 h the same pattern of letters applied: the 0.5\% and 1\% groups were distinguished from each other, with the control group intermediate.

The raw medians are reported separately, in mg/dL, and their order differed between the two durations. At 24 h the 1\% group had the highest median at 198.00, against 118.00 in the control group and 113.00 in the 0.5\% group. At 96 h the 0.5\% group had the highest median at 317.00, the control group 182.00, and the 1\% group the lowest at 155.50 (Figure~8; Tables~4 and 5). The interaction-aligned letters describe contrasts on the rank scale and are not determined by the order of these medians.

Condition factor was recorded as a descriptive control measure, and its distribution across treatment and exposure groups is shown in Figure~S7.

\end{document}
